Empirical connectomes are reported to outperform degree-preserving rewirings as reservoir
computers --- to have, in the term we use throughout, positive \emph{computational
coupling}: a score gap attributable to wiring and nothing else --- and that result is
increasingly invoked as a reason to build neuromorphic hardware on connectome-derived
wiring. We find that the comparison it rests on is not one experiment but a composition of
four choices --- the null ensemble, the unit model, the operating point and the readout
protocol --- each normally reported in a subordinate clause, each able to change the sign
of the answer on identical data, and so each a part of the coupling it reports.

The first is the one we got wrong. ``Degree-preserving randomisation'' names at least two
ensembles, depending on whether an edge's weight travels with it or is re-assigned
afterwards; both preserve the degree sequence and the weight distribution exactly, and the
phrase does not distinguish them. Over $\NUnitEvaluations$
evaluations --- $\NUnitNetworks$ connectomes $\times$ $\NUnitInstances$ paired instances
$\times$ $\NUnitGains$ gains $\times$ two unit models --- the two give the connectome's
advantage opposite signs in $\NNullSignFlip$ of $\NNullCells$ cells and opposite
\emph{significant} verdicts in $\NNullVerdictFlip$ (median disagreement
$\NullDisagreeMedian$ capacity units, maximum $\NullDisagreeMax$). Against the ensemble
the published toolbox builds the connectome leads in $\NBctRateLeads$ of
$\NBctRateNetworks$ networks; against the other, on the same graphs at the same operating
point, in $\NRankRateLeads$ of $\NRankRateNetworks$. The cause is that rank-reassignment
\emph{installs} a weight--degree association the data lacks ($\NullShiftRank$ median shift
against $\NullShiftBct$), which predicts that the two must coincide on an unweighted
network --- and on the one binary network in the set they agree to $0.008$ capacity units.

We do not read this as a refutation of the published finding --- our harness differs in
other respects too --- but as evidence that a reproduction which builds its null from a
published description cannot tell which hypothesis it tested. We
report the remaining three choices on the same terms, together with a motif-preserving
null ensemble, a gap-free perturbation bound whose prediction our own test partly
falsifies, and a full audit of the withdrawn claims of an earlier version of this work.
